\section{Gaussian transport and tensor derivatives}
\label{sec:jets}

The next step toward Theorem~\ref{thm:primitive} is to describe a
transport from a Gaussian to the auxiliary law. Its velocity will be a
conditional gradient mean. Although the original potential has only two
derivatives, conditioning by a nondegenerate Gaussian makes that mean
smooth. We will bound each of its derivative tensors as a matrix under
every division of its slots into two groups. These bounds keep later
tensor contractions from introducing a dimension factor at each step.
The interpolation itself is the Gaussian stochastic interpolant of
\citet[Section~5.1, Equations~(5.5)--(5.6)]{albergo2025interpolants};
we use its conditional velocity and prove the tensor estimates required
for our sampling construction.

Retain the primitive notation from Section~\ref{sec:preliminaries}, and assume
$l\le 1/2$, as holds for all sufficiently large $d$.  For independent
$Z\sim\nu$ and $N\sim N(0,I_d)$, put
\[
 a=1-\rho^2,\qquad Y=\rho Z+\sqrt a\,N,\qquad
 \xi_\rho=\mathcal L(Y),\qquad 0\le\rho<1.
\]
Write $\xi_\rho(dy)=e^{-H_\rho(y)}\,dy$, incorporating the normalizing
constant in $H_\rho$, and define
\[
 M_\rho(y)=\E[f(Z)\mid Y=y].
\]
We use the version defined at every $y$ by the conditional density.
Completing its quadratic gives
\begin{equation}
 \mathcal L(Z\mid Y=y)(dz)
 \ \propto\
 \exp\left(-\frac{|z-\rho y|^2}{2a}-F(x+rz)\right)dz.
 \label{eq:conditional-tilt}
\end{equation}
Under the change of variables $z=\rho y+\sqrt a\,w$, $M_\rho(y)$ is
exactly the primitive mean with center $x+r\rho y$ and scale $r\sqrt a$.
In particular, $M_0=u$.
Thus the conditional mean is a new instance of the problem already
defined in \eqref{eq:primitive-law}, at a specified center and a smaller
positive scale. This is a density identity; evaluating that mean from
finitely many gradient queries remains part of the construction.

\subsection{Cumulants from the Poincar\'e inequality}

For related derivative and cumulant estimates for Gaussian-smoothed
log-concave potentials, see \citet[Appendix~C.2]{chen2026picard}.
Here we need bounds for every nontrivial unfolding of each tensor;
we derive that form directly from Poincar\'e.

A tensor $T\in(\R^D)^{\otimes n}$ can be viewed as a matrix by
partitioning its $n$ slots into two nonempty sets $A$ and $B$.
Contraction on the $A$ slots then defines a linear map
$(\R^D)^{\otimes |A|}\to(\R^D)^{\otimes |B|}$, where both tensor
spaces carry their Euclidean, or Frobenius, norm. We call the operator
norm of this matrix a \emph{split norm}. Equivalently, it is the
supremum of $|T:(V\otimes W)|$ over coefficient tensors $V,W$ of
Frobenius norm at most one, with slots in $A,B$, respectively.
The colon denotes contraction over all paired slots.
A bound on $|T|_{\mathrm{split}}$ means that bound for every such
partition. Permuting slots simply permutes these norms. The Frobenius
norm of a tensor itself is denoted by $|T|_F$.

We use the multivariate Appell polynomials described in
\citet[Section~2.3]{anshelevich2004appell}. Their vanishing expected
lower derivatives connect Poincar\'e inequalities to tensor covariance
bounds; see \citet[Proposition~3.2 and Theorem~3.3]{adamczak2015concentration}
and \citet[Section~4.2]{kothari2017agnostic}. We give the direct
$L^2$ argument, with its dependence on the polynomial degree, and then
use a logarithm with two groups of variables to bound every split of
the cumulant tensor.

We use the following consequence of Poincar\'e, including for random vectors
whose laws need not have densities in their ambient spaces.

\Needspace{13\baselineskip}
\begin{lemma}[All-split cumulant bound]
\label{lem:appell-cumulants}
Let $S$ be a random vector in $\R^D$ whose moment-generating function is
finite in a neighborhood of zero. Suppose that, for some $\beta\ge0$,
\[
 \Var(P(S))\le\beta\E|\nabla P(S)|^2
\]
for every polynomial $P$ on $\R^D$. For each integer $n\ge2$, define
its cumulant tensor by
$\kappa_n(S)=\left.\nabla_\theta^n\log\E e^{\theta\cdot S}\right|_{\theta=0}$.
Then, for an absolute constant $C$,
\[
 |\kappa_n(S)|_{\mathrm{split}}
 \le C^n n^{Cn}\beta^{n/2}
\]
for every nontrivial split.  The constant is independent of $D$ and of the
mean of $S$.
\end{lemma}

\begin{proof}
We first prove a dimension-independent covariance bound for centered
polynomial tensors. We then recover cumulants from those covariances.
If $\beta=0$, applying the hypothesis to coordinate functions shows
that $S$ is almost surely constant, so every cumulant in the statement
vanishes. We may therefore suppose $\beta>0$.

Define the symmetric Appell tensor polynomials by
\[
 P_j(s)=
 \left.\nabla_\theta^j
 \frac{\exp(\theta\cdot s)}{\E\exp(\theta\cdot S)}
 \right|_{\theta=0}.
\]
For a coefficient tensor $T$ of order $j$, let $Q_T(s)=T:P_j(s)$,
where the colon denotes contraction on all slots.  Differentiating the
generating function in $s$ multiplies it by copies of $\theta$.
Consequently
\[
 \E\nabla^h Q_T(S)=0\quad(0\le h<j),\qquad
 \nabla^j Q_T=j!\operatorname{Sym}(T).
\]
Here $\operatorname{Sym}(T)$ is the average over all permutations of
the slots; hence $|\operatorname{Sym}(T)|_F\le|T|_F$.
Apply Poincar\'e to each component of $\nabla^h Q_T$ and sum.  Iterating
from $h=0$ to $h=j-1$ yields
\begin{equation}
 \E|T:P_j(S)|^2\le \beta^j(j!)^2|T|_F^2.
 \label{eq:appell-l2}
\end{equation}
Only first derivatives in the displayed Poincar\'e inequality are used at
each application.

In particular, if $T$ and $V$ have orders $j,k\ge1$, then
\begin{equation}\label{eq:appell-covariance}
 \left|\E\bigl[(T:P_j(S))(V:P_k(S))\bigr]\right|
 \le\beta^{(j+k)/2}j!k!\,|T|_F|V|_F.
\end{equation}
This is the direct covariance estimate. It concerns arbitrary
coefficient tensors, rather than products of one-dimensional test
directions, which is why it controls a matrix split.

To pass from covariance to cumulants while preserving a specified split,
consider the two-variable generating function
\[
 \mathcal G(\theta,\omega)=
 \frac{\E\exp((\theta+\omega)\cdot S)}
 {\E\exp(\theta\cdot S)\,\E\exp(\omega\cdot S)}.
\]
Its derivative of orders $j\ge1$ in $\theta$ and $k\ge1$ in $\omega$
at $(0,0)$ is the cross-moment tensor $\E[P_j(S)\otimes P_k(S)]$.
Across the split between the $\theta$ and $\omega$ slots its operator
norm is at most $\beta^{(j+k)/2}j!k!$ by
\eqref{eq:appell-covariance}.

Both $\mathcal G(\theta,0)$ and $\mathcal G(0,\omega)$ are identically
one.  Thus every nonzero term of $\mathcal G-1$ contains variables from
both groups.  In the formal identity
\[
 \log\mathcal G
 =\sum_{m\ge1}\frac{(-1)^{m+1}}m(\mathcal G-1)^m,
\]
only $m\le n$ can contribute at total order $n$.  Every differentiated
product is a tensor product of the preceding cross-moment operators,
up to permutations within the two sides.  Its norm is therefore at most
the product of their norms.  To bound the total coefficient, distribute
the $n$ labeled derivative slots among at most $n$ ordered nonempty
factors.  There are at most $n^{n+1}$ such assignments; the product of
the factorials within any assignment is at most $n!$.  Enlarging an
absolute constant bounds their total by $C^n n^{Cn}$.

The mixed derivatives of $\log\mathcal G$ are precisely the
corresponding derivatives of
$\log\E\exp((\theta+\omega)\cdot S)$, since the other two logarithms
have only one variable group.  They are the cumulants, and the two
groups can represent any prescribed nontrivial split of the $n$ slots.
This proves the claim.
\end{proof}

In the conditional density, $y$ enters through the linear tilt
$\rho y/a$, so each differentiation of $M_\rho$ with respect to $y$
supplies a factor $\rho$. We divide out these factors to state bounds
that remain meaningful at $\rho=0$.

\Needspace{19\baselineskip}
\begin{proposition}[Gaussian jet bounds]
\label{prop:gaussian-jets}
Under the assumptions above, in particular $l\le1/2$, define for
each integer $q\ge1$
\[
 U_q(\rho,y)=\rho^{-q}\nabla^q M_\rho(y)\qquad(\rho>0).
\]
These tensors extend smoothly to $\rho=0$. There is an absolute
constant $C$ such that, for every $0\le\rho<1$, $y\in\R^d$ and
every split,
\begin{equation}
 |U_q(\rho,y)|_{\mathrm{split}}
 \le C^q q^{Cq}L(\sqrt a)^{-(q-1)}.
 \label{eq:jets}
\end{equation}
Here and below the vector-output slot is included among the tensor slots.
The constant is uniform in the dimension, the primitive, its center and
scale, and the derivative order $q$.
Moreover, $DM_\rho$ and $U_1$ are symmetric, and
\begin{equation}
 \nabla H_\rho(y)=y+r\rho M_\rho(y),\qquad
 \|\nabla^2H_\rho\|_{\op}\le C.
 \label{eq:interpolation-score}
\end{equation}
For $j\ge1$, every split of $\nabla^j\nabla^2 H_\rho$ is bounded by
\begin{equation}
 C^{j+1}(j+1)^{C(j+1)}(\sqrt a)^{-j}.
 \label{eq:score-derivatives}
\end{equation}
\end{proposition}

\begin{proof}
The conditional potential in \eqref{eq:conditional-tilt} has Hessian
between $(a^{-1}-l)I$ and $(a^{-1}+l)I$.  By
Lemma~\ref{lem:convex-preliminaries}, its Poincar\'e constant is at most
$Ca$.  Under this conditional law the map
\[
 z\longmapsto S(z)=(z,f(z)/L)
\]
has derivative norm at most $\sqrt2$, since $f$ is $L$-Lipschitz.
Pulling back an ambient test function shows that the possibly singular
law of $S$ satisfies Poincar\'e with constant $Ca$.  It has exponential
moments of every linear functional: $f$ has at most linear growth and
the conditional law has Gaussian tails.  Lemma~\ref{lem:appell-cumulants}
therefore applies.  This use of repeated Poincar\'e does not differentiate
$f$ repeatedly.  Each application pulls back a new ambient polynomial
and uses only $Df$.

In the conditional density, the linear tilt parameter for $z$ is
$\rho y/a$.  Repeated differentiation of an expectation with respect
to a linear tilt gives joint cumulants.  In coordinates this says
\begin{equation}
 (U_q)_{i,k_1,\ldots,k_q}
 =a^{-q}\operatorname{cum}
       (f_i(Z),Z_{k_1},\ldots,Z_{k_q}\mid Y=y).
 \label{eq:jet-cumulant}
\end{equation}
It can also be checked directly by differentiating the logarithm of the
conditional moment-generating function with one auxiliary parameter
coupled to $f_i$.  Restrict the cumulant of $S$ to one $f/L$ coordinate
and $q$ ordinary coordinates.  Lemma~\ref{lem:appell-cumulants} gives a
factor $L a^{(q+1)/2}$, and the prefactor $a^{-q}$ gives
\eqref{eq:jets}.

All the conditional integrals used here are smooth in $(\rho,y)$ for
$\rho<1$.  This follows by differentiating their densities: on compact
parameter sets, every derivative inserts a polynomial with an integrable
Gaussian majorant.  Formula~\eqref{eq:jet-cumulant} gives the smooth
extension at zero without division by $\rho$.
For later integrations by parts we record a slightly stronger
regularity fact.  On each compact interval in $\rho<1$, conditional
moments of any fixed order grow at most polynomially in $y$.  Indeed,
the conditional mode satisfies
\[
 z_0=\rho y-ar g(x+rz_0),
\]
so $|z_0|\le C(|y|+r|g(x)|)$.  The upper and lower conditional Hessian
bounds compare the density and its normalizing integral around this
mode to Gaussian integrals.  This proves the claimed moment growth.
The same argument applies to time derivatives of the conditional
integrals.  Dimension-dependent constants are harmless for this
regularity justification; the quantitative estimates above remain
dimension-uniform.

Conditional integration by parts gives
$\E[Z\mid Y=y]=\rho y-arM_\rho(y)$.  Differentiating the marginal
density then yields the first identity in
\eqref{eq:interpolation-score}.  For $\rho>0$ it implies symmetry of
$DM_\rho$, hence of $U_1$; continuity handles zero.  Its derivative,
together with \eqref{eq:jets}, gives the bounded Hessian and
\eqref{eq:score-derivatives}.  In particular, all these conclusions use
only the original $C^2$ regularity of $F$.
\end{proof}

\subsection{The transport equation}

\begin{proposition}[Probability transport]
\label{prop:transport}
The ordinary differential equation
\begin{equation}
 \frac{dY_\rho}{d\rho}=-rM_\rho(Y_\rho)
 \label{eq:transport}
\end{equation}
transports $\xi_0=N(0,I_d)$ to $\xi_\rho$ for every $\rho<1$.
Its velocity is globally $Cl$-Lipschitz in space.  On every compact
probability-time interval it has a smooth flow in both directions.
\end{proposition}

\begin{proof}
For the random interpolation used to define $\xi_\rho$,
\[
 \frac{d}{d\rho}(\rho Z+\sqrt a\,N)
 =\frac{Z-\rho Y}{a}.
\]
The conditional mean of this derivative is $-rM_\rho(Y)$ by the
conditional integration-by-parts identity above.  Thus, for compactly
supported smooth $\phi$,
\[
 \frac{d}{d\rho}\int\phi\,d\xi_\rho
 =-r\int M_\rho\cdot\nabla\phi\,d\xi_\rho.
\]
On the other hand, \eqref{eq:jets} with $q=1$ gives
$r\|DM_\rho\|_{\op}\le Cl$.  The velocity has at most linear growth,
so its smooth ordinary differential equation has a global-in-space
flow and inverse on every compact interval in $\rho<1$.
To identify its pushforward, fix a smooth compactly supported terminal
test and transport it backward along this flow.  The transported test
has zero material derivative.  Its supports remain in a bounded set
on the interval, by the linear-growth bound.  The last display then
shows its expectation under $\xi_\rho$ is constant, which proves the
pushforward assertion.
\end{proof}

The remaining analytic difficulty is not spatial smoothness, which the
Gaussian interpolation has supplied.  It is to control repeated
derivatives \emph{along} the transport while paying only one
$\sqrt d$ factor for vector outputs.  We establish that control next.
